\documentclass[11pt]{article}

\usepackage[T1]{fontenc}
\usepackage[margin=1in]{geometry}
\usepackage{amsmath, amssymb, amsthm}
\usepackage{enumitem}
\usepackage{hyperref}
\usepackage{booktabs}
\usepackage{listings}
\usepackage{xcolor}

\lstset{
  basicstyle=\ttfamily\small,
  frame=single,
  breaklines=true,
  columns=fullflexible,
  mathescape=true
}

\theoremstyle{definition}
\newtheorem{theorem}{Theorem}[section]
\newtheorem{definition}[theorem]{Definition}
\newtheorem{example}[theorem]{Example}

\title{CS 250 --- Module 6: Modular Arithmetic and Proofs by Cases}
\author{}
\date{}

\begin{document}
\maketitle

\section{The modulus operation}
An operation that's important in both mathematics and computer science is the modulus\index{modulus} operator, sometimes called \texttt{mod} or \texttt{\%}. This is the operation that gives you the remainder after integer division.

Examples:
\begin{lstlisting}
5 % 2 = 1
9 % 3 = 0
59 % 4 = 3
\end{lstlisting}

The modulus operation takes the natural numbers---which are NOT a group under addition, because there are no additive inverses---and gives them really nice properties!

\section{The quotient-remainder theorem}

Let's make this precise.

\begin{theorem}[Quotient-Remainder Theorem]\index{quotient-remainder theorem}
For any integer $n$ and any positive integer $d$, there exist \textbf{unique} integers $q$ and $r$ such that
\[
  n = d \cdot q + r \quad \text{and} \quad 0 \le r < d.
\]
We call $q$ the \emph{quotient} (this is integer division) and $r$ the \emph{remainder} (this is the mod operation).
\end{theorem}

This theorem formalizes the quotient-remainder decomposition: the decomposition into quotient and remainder always exists and is always unique.

If you're a programmer you already know this in your bones. In Python, \texttt{//} gives you $q$ and \texttt{\%} gives you $r$. In C, \texttt{/} on ints gives you $q$ and \texttt{\%} gives you $r$---at least for non-negative values. For negative numbers, C's division truncates toward zero rather than rounding down, so the remainder can be negative; Python's \texttt{//} and \texttt{\%} always match the theorem. So the quotient-remainder theorem is literally just this:
\begin{lstlisting}
n == d * (n // d) + (n % d)    # always True!
\end{lstlisting}
That's it. That's the theorem, in code.

Let's do a worked example. Take $n = 67$ and $d = 4$. Then $67 = 4 \cdot 16 + 3$, so $q = 16$ and $r = 3$. You can check: $4 \times 16 = 64$, and $64 + 3 = 67$. And $0 \le 3 < 4$, so the remainder condition is satisfied.

Why does this matter? The theorem guarantees that every integer falls into exactly one \emph{residue class} mod $d$: the remainder is always one of $0, 1, 2, \ldots, d-1$. This is why proof by cases using mod works! For example, every integer is either even ($r = 0$ mod 2) or odd ($r = 1$ mod 2)---there's no third option. The quotient-remainder theorem is what makes that airtight.

\section{Floor and ceiling}

Two functions that show up constantly in both math and programming are \emph{floor} and \emph{ceiling}.

\begin{definition}[Floor and Ceiling]
The \emph{floor}\index{floor} of a real number $x$, written $\lfloor x \rfloor$, is the largest integer that is less than or equal to $x$---``rounding down.'' In Python this is \texttt{math.floor(x)}, or for positive numbers just \texttt{int(x)}.

The \emph{ceiling}\index{ceiling} of $x$, written $\lceil x \rceil$, is the smallest integer that is greater than or equal to $x$---``rounding up.'' In Python: \texttt{math.ceil(x)}.
\end{definition}

Some examples:
\[
  \lfloor 3.7 \rfloor = 3, \quad \lceil 3.7 \rceil = 4, \quad \lfloor 5 \rfloor = 5, \quad \lceil 5 \rceil = 5, \quad \lfloor -2.3 \rfloor = -3, \quad \lceil -2.3 \rceil = -2.
\]

One key property: $\lfloor x \rfloor = x$ if and only if $x$ is an integer.\footnote{This sounds obvious but you'll need to prove it carefully on the assignment!} If $x$ has any fractional part at all, the floor will be strictly less than $x$.

There's a nice connection to the quotient-remainder theorem too. For positive integers $n$ and $d$, the quotient $q$ is exactly $\lfloor n / d \rfloor$. So floor and integer division are really the same thing.

Fun fact: the floor function shows up all over computer science. Array indexing, hash functions, binary search---anywhere you need to convert a real-valued quantity into an integer index, you're probably using floor.

\bigskip
\textbf{Exercise 1.} \label{ex:m6:1} Compute $\lfloor -3.7 \rfloor$, $\lceil -3.7 \rceil$, $\lfloor 4 \rfloor$, and $\lceil 4 \rceil$. (Be careful with the negative value---remember that floor rounds \emph{down}, not toward zero.)

\section{$\mathbb{Z}/n\mathbb{Z}$ is a group}
In Module 5 we proved that $(\mathbb{Z}, +, 0)$ is a group. Now we'll see that modular arithmetic gives us \emph{finite} groups---same axioms, but with a finite set of residues instead of all the integers.

For any positive integer $n$, we can define the set of residues $\mathbb{Z}/n\mathbb{Z} = \{0, 1, \ldots, n-1\}$. This set \textbf{is} actually a group under addition. We'll write $\oplus$ for the modular addition to distinguish it from ordinary $+$:
\begin{lstlisting}
q, r : Z % n
q (+) r = (q + r) % n
\end{lstlisting}

In other words, we perform ordinary addition and then take the modulus. (In practice, most texts just write $+$ for both and let context disambiguate, but the distinction is worth noting.)

The group axioms work the way you'd hope: associativity comes from ordinary integer addition, the identity is still 0, and every residue has an additive inverse.

For every $q \in \mathbb{Z}/n\mathbb{Z}$, the inverse is the element $r = (n - q) \bmod n$.

Why is this the inverse? Because
\begin{lstlisting}
q (+) (n - q) = (q + n - q) % n
              = n % n
              = 0
\end{lstlisting}

\bigskip
\textbf{Exercise 2.} \label{ex:m6:2} In $\mathbb{Z}/7\mathbb{Z}$, find the additive inverse of $3$. Verify your answer by checking that the sum of $3$ and your proposed inverse is $0$ under modular addition.

\section{$\mathbb{Z}/p\mathbb{Z}$, where $p$ is prime}
When the number you're taking the modulus by is prime, something \textbf{really} special happens. It isn't just a group under addition---if you exclude 0, the nonzero residues also form a group under multiplication. That is, $((\mathbb{Z}/p\mathbb{Z})^{\times}, \cdot, 1)$ satisfies exactly the same group axioms we checked for the integers in Module 5, just with a different set and a different operation.
\begin{lstlisting}
q, r : Z % p, q /= 0, r /= 0
q * r = (q * r) % p
\end{lstlisting}
\begin{theorem}[B\'ezout's Identity]\index{Bezout's identity@B\'ezout's identity}
For any integers $a$ and $b$, there exist integers $x$ and $y$ such that $ax + by = \gcd(a, b)$.
\end{theorem}
We won't prove this here, but Module 7 shows how the extended Euclidean algorithm computes these coefficients.

It's maybe not obvious that multiplicative inverses exist, but they do: if $q$ is a nonzero residue mod $p$ and $p$ is prime, then $\gcd(q,p) = 1$, so B\'ezout's identity gives integers $x,y$ with $qx + py = 1$. Reducing mod $p$ then shows $qx \equiv 1 \pmod p$, so $x$ is a multiplicative inverse of $q$ modulo $p$.

Distributivity of multiplication over addition follows from distributivity in $\mathbb{Z}$, so we won't re-prove it here.

The importance of this is that it makes $\mathbb{Z}/p\mathbb{Z}$ a \emph{finite field}\index{finite field}: addition works on all residues, and multiplication works on the nonzero residues. This is actually really important in cryptography!

\bigskip
\textbf{Exercise 3.} \label{ex:m6:3} In $\mathbb{Z}/5\mathbb{Z}$, find the multiplicative inverse of $3$. Verify your answer by checking that $3 \cdot x \equiv 1 \pmod{5}$.

For further reading on finite fields and their applications:

\url{https://en.wikipedia.org/wiki/Finite_field_arithmetic}

\url{https://en.wikipedia.org/wiki/Finite_field#Discrete_logarithm}

\section{Proofs by cases and the connection to or-elimination}
Proof by cases is $\lor$-elimination in disguise. If you have an exhaustive set of cases for the data of your problem---like a number is either 0 or non-zero, or a number is positive, zero, or negative, or a number is either odd or even, \&c.---then you can treat the disjunction of those cases as a tautology and apply $\lor$-elimination to it.

Let's do a concrete example. We'll prove that for every integer $n$, the quantity $n^2 - n + 3$ is odd. The quotient-remainder theorem tells us every integer is either even or odd (the two residue classes mod 2), so those are our two cases.

\textbf{Case 1:} $n$ is even, so $n = 2k$ for some integer $k$. Then
\[
  n^2 - n + 3 = (2k)^2 - 2k + 3 = 4k^2 - 2k + 3 = 2(2k^2 - k + 1) + 1,
\]
which is odd.

\textbf{Case 2:} $n$ is odd, so $n = 2k + 1$ for some integer $k$. Then
\[
  n^2 - n + 3 = (2k+1)^2 - (2k+1) + 3 = 4k^2 + 4k + 1 - 2k - 1 + 3 = 4k^2 + 2k + 3 = 2(2k^2 + k + 1) + 1,
\]
which is odd.

So in both cases we get something of the form $2m + 1$, which is odd. Since every integer is either even or odd, we're done. Notice how the quotient-remainder theorem is doing the heavy lifting here---it's the reason we know these two cases are \emph{exhaustive}. We don't have to worry about some integer that's neither even nor odd, because the theorem guarantees there isn't one.

\end{document}
